\documentclass[10pt,a4paper]{amsart}
\usepackage[margin=19mm]{geometry}
\usepackage{amsmath,amssymb,mathtools}
\usepackage[scr=rsfs,cal=euler]{mathalfa}
\usepackage{xfrac}
\usepackage[unicode,hidelinks,bookmarks=false]{hyperref}
\newcommand{\ie}{i.e.\ }
\newcommand{\cf}{cf.\ }
\newcommand{\resp}{respectively\ }
\newcommand{\Subsection}{\subsection}
\providecommand{\nobreakdash}{\nobreak\hbox{-}}
\setlength{\emergencystretch}{2em}
\multlinegap0pt
\usepackage{siunitx}
\usepackage{siunitx}
\usepackage{siunitx}
\usepackage{siunitx}
\usepackage{bm}
\usepackage{mathtools}
\usepackage{algorithm}\usepackage{algpseudocode}
\usepackage[shortlabels]{enumitem}
\usepackage{siunitx}
\usepackage{xcolor}
\newtheorem{theorem}{Theorem}
\title{Unleashing LLMs in Bayesian Optimization: Preference-Guided Framework for Scientific Discovery}
\author{Xinzhe Yuan, Zhuo Chen, Jianshu Zhang, Huan Xiong, Nanyang Ye, Yuqiang Li, Qinying Gu}
\date{Source: arXiv:2605.17976v1 (CC BY 4.0)}
\begin{document}
\maketitle

\noindent\emph{Source and licence.} Unleashing LLMs in Bayesian Optimization: Preference-Guided Framework for Scientific Discovery. Version arXiv:2605.17976v1 (CC BY 4.0), distributed by arXiv under the Creative Commons Attribution 4.0 International licence, http://creativecommons.org/licenses/by/4.0/, as recorded in the arXiv licence metadata for this exact version and shown on the abstract page https://arxiv.org/abs/2605.17976v1. Full licence notice: \texttt{licenses/arxiv-2605.17976-CC-BY-4.0.txt}.\par\smallskip

\noindent\textit{Editorial scope.} Counted unit: the article's theorem app:the2 (source0.tex lines 1108--1128). The article's complete original proof of it is delivered with it (source0.tex lines 1130--1157, one \texttt{proof} environment of 28 lines). No mathematics is written, repaired or re-derived: the statement and the proof are the article's own lines, in the article's own notation, and no retained line is substituted or re-flowed.  No further source-local context is needed: every cross-reference of every retained line resolves inside the two delivered spans.  Nothing was omitted from any retained span, so the package carries no omission record.  No retained line contains a citation, so the wrapper carries no bibliography and no bibliography entry.  No retained line contains a figure environment, an image-inclusion command or drawing code, so no image is delivered and the package declares no asset dependency.  Editorial formatting only, in addition: an AMS \texttt{amsart} preamble that restates the article's own declarations and macros used by the retained lines, namely the environments \texttt{theorem} on separate counters, together with the standard packages those lines need.  The retained environments print in this document as Theorem 1 in order of appearance, whereas the article prints the same items with the numbering of its own full document.  The complete native source of the article as distributed in the arXiv e-print for this exact version is bundled unchanged as \texttt{sources/200828/source0.tex}.
\begin{theorem}[Gaussian measure under exponential linear lift]\label{app:the2}
 Let $(E,\mathcal B)$ be a separable real topological vector space, $E^*$ its continuous dual, and $\mathbb P=\mathcal N(m,C)$ a Gaussian measure with mean $m\in E$ and covariance operator $C:E^*\to E$.  
For any continuous linear functional $\ell\in E^*$ and any $\lambda\in\mathbb R$, define the lifted measure
\[
\frac{d\mathbb P_\lambda}{d\mathbb P}(f)
= \exp\!\big(\lambda\,\ell(f)-\psi(\lambda)\big),\qquad
\psi(\lambda)=\log\mathbb E_\mathbb P[e^{\lambda \ell(f)}].
\]

Then $f|\mathbb P_\lambda$ is still Gaussian with unchanged covariance $C$, and mean shifted by
\[
m_\lambda = m + \lambda C\ell.
\]
Moreover, the expected lift of the functional is
\[
\Delta = \mathbb E_{\mathbb P_\lambda}[\ell(f)] - \mathbb E_\mathbb P[\ell(f)]
= \lambda\,\langle C\ell,\ell\rangle
= \lambda\,\mathrm{Var}_\mathbb P[\ell(f)].
\]
Thus a prescribed increment $\delta$ is achievable iff $\mathrm{Var}_\mathbb P[\ell(f)]>0$, in which case $\lambda=\Delta/\mathrm{Var}_\mathbb P[\ell(f)]$.
\end{theorem}

\begin{proof}
We present the argument via finite-dimensional projections.

For any finite set of functionals $\Phi=(\phi_1,\dots,\phi_m)\subset E^*$, the vector
\[
Y=(\phi_1(f),\dots,\phi_m(f),\,\ell(f))
\]
is multivariate Gaussian under $\mathbb P$.  
Lifting the density by $\exp(\lambda \ell(f))$ yields, by standard Gaussian completion-of-squares or cumulant-generating function arguments, that
\[
(\phi_1(f),\dots,\phi_m(f))^\top \;\sim\; 
\mathcal N\!\big( \mu_\Phi + \lambda\,\mathrm{Cov}_\mathbb P((\phi_i(f))_i,\ell(f)),\ \Sigma_\Phi \big).
\]
Here $\mu_\Phi=(\phi_i(m))_i$ and $\Sigma_\Phi=(\langle C\phi_i,\phi_j\rangle)_{ij}$.  
But
\[
\mathrm{Cov}_\mathbb P(\phi(f),\ell(f))=\langle C\ell,\phi\rangle,
\]
so the new mean is exactly $(\phi_i(m+\lambda C\ell))_i$, while the covariance remains $\Sigma_\Phi$.  
By Kolmogorov consistency, this determines a Gaussian measure $\mathcal N(m+\lambda C\ell, C)$ on $E$.  

Finally,
\[
\Delta=\ell(m+\lambda C\ell)-\ell(m)
=\lambda\,\langle C\ell,\ell\rangle,
\]
which equals $\lambda\,\mathrm{Var}_\mathbb P[\ell(f)]$, yielding the desired characterization of $\lambda$.
\end{proof}

\end{document}
